\noindent\textbf{Elección del procedimiento.} Para una etapa ideal sin fuerzas exteriores, $dv=-c\,dm/m$ con velocidad efectiva constante $c=I_{sp}g_0$. Integrar de $m_0$ a $m_f$ da $\Delta v=c\ln(m_0/m_f)$.

\begin{enumerate}

\item Convertir impulso específico a velocidad efectiva.

\[\fitmath{c=\mathrm{Isp}\cdot g_{0}}\]

\[\fitmath{c=(300)\cdot (9.80665)\approx 2941.995\,\mathrm{m/s}}\]

\item El argumento del logaritmo debe ser adimensional y mayor que uno.

\[\fitmath{\mathrm{ratio}=\frac{m_{0}}{\mathrm{mf}}}\]

\[\fitmath{\mathrm{ratio}=\frac{(500000)}{(125000)}\approx 4\,}\]

\item La integral de $1/m$ es logaritmo natural, no decimal.

\[\fitmath{\mathrm{dv}=c\cdot \ln\left(\mathrm{ratio}\right)}\]

\[\fitmath{\mathrm{dv}=(2941.995)\cdot \ln\left((4)\right)\approx 4078.471\,\mathrm{m/s}}\]

\end{enumerate}
